\documentclass{article}
\usepackage{amsmath,amssymb,amsthm}
\usepackage{algorithm}
\usepackage[noend]{algpseudocode}
\usepackage{enumitem}
\newtheorem{theorem}{Theorem}
\newtheorem{lemma}{Lemma}
\newtheorem{assumption}{Assumption}
\newcommand{\norm}[1]{\left\lVert#1\right\rVert}
\begin{document}

\section*{Algorithm Name}
\textbf{Trust‑Region Gradient Method (TR‑GM)} for convex and smooth objectives.

\section*{Mathematical Formulation}
Let $f:\mathbb{R}^{d}\rightarrow\mathbb{R}$ be the objective.  
At iteration $k$ we build a quadratic model
\[
m_k(p)=f(x_k)+\nabla f(x_k)^{\top}p+\tfrac12 p^{\top}B_k p,
\qquad p\in\mathbb{R}^{d},
\]
where $B_k$ is a symmetric positive‑semidefinite matrix (typically $B_k=\nabla^{2}f(x_k)$ or $B_k=L I$).  
The trial step $p_k$ is obtained by (approximately) solving the trust‑region subproblem
\[
\min_{p}\; m_k(p)\quad\text{s.t.}\quad \norm{p}\le \Delta_k,
\tag{TR}
\]
with trust‑region radius $\Delta_k>0$.

Define the \emph{actual reduction}
\[
\operatorname{ared}_k = f(x_k)-f(x_k+p_k),
\]
and the \emph{predicted reduction}
\[
\operatorname{pred}_k = m_k(0)-m_k(p_k).
\]
The ratio
\[
\rho_k = \frac{\operatorname{ared}_k}{\operatorname{pred}_k}
\tag{1}
\]
governs acceptance of the trial step and update of $\Delta_k$:
\[
x_{k+1}= 
\begin{cases}
x_k+p_k, & \rho_k\ge \eta_1,\\[4pt]
x_k,    & \rho_k<\eta_1,
\end{cases}
\tag{2}
\]
\[
\Delta_{k+1}= 
\begin{cases}
\min\{\gamma_{\text{inc}}\Delta_k,\;\Delta_{\max}\}, & \rho_k\ge \eta_2,\\[4pt]
\gamma_{\text{dec}}\Delta_k, & \rho_k<\eta_1,
\end{cases}
\tag{3}
\]
with parameters $0<\eta_1<\eta_2<1$, $0<\gamma_{\text{dec}}<1<\gamma_{\text{inc}}$, and $\Delta_{\max}>0$.

\section*{Pseudocode}
\begin{algorithm}[H]
\caption{Trust‑Region Gradient Method (TR‑GM)}
\begin{algorithmic}[1]
\State \textbf{Input:} $x_0$, initial radius $\Delta_0>0$, parameters
$\eta_1,\eta_2,\gamma_{\text{dec}},\gamma_{\text{inc}},\Delta_{\max}$.
\For{$k=0,1,2,\dots$}
    \State Compute $\nabla f(x_k)$.
    \State Choose $B_k\succeq 0$ (e.g., $B_k=L I$).
    \State Approximately solve \eqref{TR} to obtain $p_k$.
    \State Compute $\rho_k$ via \eqref{1}.
    \If{$\rho_k\ge \eta_1$}
        \State $x_{k+1}\gets x_k+p_k$.
    \Else
        \State $x_{k+1}\gets x_k$.
    \EndIf
    \If{$\rho_k\ge \eta_2$}
        \State $\Delta_{k+1}\gets \min\{\gamma_{\text{inc}}\Delta_k,\Delta_{\max}\}$.
    \ElsIf{$\rho_k<\eta_1$}
        \State $\Delta_{k+1}\gets \gamma_{\text{dec}}\Delta_k$.
    \Else
        \State $\Delta_{k+1}\gets \Delta_k$.
    \EndIf
    \If{termination criterion satisfied (e.g., $\norm{\nabla f(x_{k+1})}\le\varepsilon$)}
        \State \textbf{break}
    \EndIf
\EndFor
\State \textbf{Output:} $x_{k+1}$.
\end{algorithmic}
\end{algorithm}

\section*{Dry Run Example}
Consider the one‑dimensional convex quadratic
\[
f(w)=(w-3)^2,\qquad w\in\mathbb{R},
\]
with $L=2$ (Lipschitz constant of $\nabla f$).  
Take $w_0=0$, initial radius $\Delta_0=1$, and parameters
$\eta_1=0.25,\;\eta_2=0.75,\;\gamma_{\text{dec}}=0.5,\;\gamma_{\text{inc}}=2,\;\Delta_{\max}=4$.
We use $B_k = L =2$ (scalar) and solve the subproblem analytically:
\[
p_k = \arg\min_{|p|\le\Delta_k}\;\bigl[(w_k-3)^2 + 2(w_k-3)p + p^2\bigr]
      =\operatorname{proj}_{[-\Delta_k,\Delta_k]}\bigl(-\nabla f(w_k)\bigr).
\]

\begin{enumerate}[label=\textbf{Iter \arabic*:},leftmargin=*]
\item $k=0$  
\[
\nabla f(w_0)=2(w_0-3)=-6,\qquad p_0=\operatorname{proj}_{[-1,1]}(6)=1.
\]
Model prediction:
\[
\operatorname{pred}_0 = m_0(0)-m_0(p_0)
 = f(w_0)-\bigl[f(w_0)+\nabla f(w_0)p_0+\tfrac12 B_0p_0^2\bigr]
 = -\bigl[-6\cdot1+\tfrac12\cdot2\cdot1^2\bigr]=5.
\]
Actual reduction:
\[
f(w_0)=9,\;f(w_0+p_0)=f(1)=4,\qquad 
\operatorname{ared}_0 =9-4=5.
\]
Thus $\rho_0=5/5=1\ge\eta_2$, accept step and enlarge radius:
\[
w_1=0+1=1,\qquad \Delta_1=\min\{2\cdot1,4\}=2.
\]

\item $k=1$  
\[
\nabla f(w_1)=2(1-3)=-4,\qquad p_1=\operatorname{proj}_{[-2,2]}(4)=2.
\]
\[
\operatorname{pred}_1 = -\bigl[-4\cdot2+\tfrac12\cdot2\cdot2^2\bigr]=4.
\]
\[
f(w_1)=4,\;f(w_1+p_1)=f(3)=0,\qquad 
\operatorname{ared}_1 =4-0=4.
\]
$\rho_1=4/4=1\ge\eta_2$, accept step, enlarge radius:
\[
w_2=1+2=3,\qquad \Delta_2=\min\{2\cdot2,4\}=4.
\]

\item $k=2$  
\[
\nabla f(w_2)=2(3-3)=0,\qquad p_2=0.
\]
Both predicted and actual reductions are $0$, $\rho_2$ is undefined; we treat it as $1$ and keep the iterate:
\[
w_3=w_2=3,\qquad \Delta_3=\min\{2\cdot4,4\}=4.
\]
The algorithm has reached the minimizer $w^\star=3$ in two effective steps.
\end{enumerate}

\section*{Assumptions}
\begin{assumption}[Convexity and Smoothness]\label{assump:convex}
The objective $f:\mathbb{R}^{d}\to\mathbb{R}$ is convex and differentiable, and its gradient is $L$‑Lipschitz:
\[
\|\nabla f(x)-\nabla f(y)\|\le L\|x-y\|,\qquad\forall x,y\in\mathbb{R}^{d}.
\]
Equivalently,
\[
0\preceq \nabla^{2}f(x)\preceq L I,\qquad\forall x.
\]
\end{assumption}
\begin{assumption}[Model Accuracy]\label{assump:model}
For all $k$, the matrix $B_k$ satisfies
\[
0\preceq B_k\preceq L I .
\]
Consequently, the model $m_k$ is an upper bound on $f$ within the trust region:
\[
f(x_k+p)\le m_k(p),\qquad \forall p\text{ with }\|p\|\le\Delta_k .
\]
\end{assumption}
\begin{assumption}[Subproblem Approximation]\label{assump:approx}
The computed step $p_k$ satisfies the \emph{Cauchy decrease} condition
\[
m_k(0)-m_k(p_k)\;\ge\;\kappa_{\mathrm{c}}\,
\min\!\Bigl\{\frac{\|\nabla f(x_k)\|^{2}}{L},\;
\Delta_k\|\nabla f(x_k)\|\Bigr\},
\tag{4}
\]
for some constant $\kappa_{\mathrm{c}}\in(0,1]$ (e.g., the exact Cauchy point yields $\kappa_{\mathrm{c}}=1/2$).
\end{assumption}
\begin{assumption}[Parameters]\label{assump:params}
The algorithm parameters satisfy
\[
0<\eta_1<\eta_2<1,\qquad
0<\gamma_{\text{dec}}<1<\gamma_{\text{inc}},\qquad
\Delta_{\max}>0.
\]
\end{assumption}

\section*{Main Convergence Theorem}
\begin{theorem}[Global Convergence and Rate]\label{thm:main}
Under Assumptions \ref{assump:convex}–\ref{assump:params}, the sequence $\{x_k\}$ generated by TR‑GM satisfies
\[
f(x_k)-f^\star \;\le\; \frac{C}{k+1},
\qquad\forall k\ge0,
\tag{5}
\]
where $f^\star=\min_x f(x)$ and
\[
C:=\max\Bigl\{\,\tfrac{2L}{\kappa_{\mathrm{c}}\,\eta_1}\,
\Delta_{0}^{2},\; (f(x_0)-f^\star)\Bigr\}.
\]
Moreover, $\lim_{k\to\infty}\norm{\nabla f(x_k)}=0$ and every accumulation point of $\{x_k\}$ is a minimizer of $f$.
\end{theorem}
\begin{theorem}[Complexity in Terms of Gradient Norm]\label{thm:grad}
For any tolerance $\varepsilon>0$, the algorithm produces an iterate $x_k$ with
$\norm{\nabla f(x_k)}\le\varepsilon$ after at most
\[
\mathcal{O}\!\Bigl(\frac{L\Delta_0^2}{\kappa_{\mathrm{c}}\eta_1\varepsilon^{2}}\Bigr)
\]
successful iterations.
\end{theorem}

\section*{Theoretical Properties}
\begin{itemize}
\item \textbf{Stability.}  Because $B_k\preceq L I$, the model never over‑estimates curvature; the trust‑region mechanism guarantees that the step size never exceeds the radius, preventing divergence even when the gradient is large.
\item \textbf{Hyper‑parameter Sensitivity.}  The constants $\eta_1,\eta_2$ only affect the frequency of radius updates; any choice satisfying $0<\eta_1<\eta_2<1$ yields the same $O(1/k)$ rate, although larger $\eta_2$ typically reduces the number of radius expansions.
\item \textbf{Comparison to Gradient Descent.}  Setting $\Delta_k=\infty$ reduces TR‑GM to exact Newton (or gradient) steps, reproducing the standard $O(1/k)$ rate for convex $L$‑smooth functions. The trust region, however, adds robustness when $B_k$ is only an approximation.
\end{itemize}

\section*{Discussion}
The analysis shows that, despite solving the subproblem only approximately, the Cauchy‑type guarantee \eqref{4} together with convexity suffices to obtain the classical $1/k$ decay of the optimality gap. The result matches the optimal rate for first‑order methods on convex smooth problems, while providing a safeguard against poor curvature approximations.

\newpage
\section*{Preliminaries (Definitions and Lemmas)}
\begin{lemma}[Descent Lemma]\label{lem:descent}
If $f$ is $L$‑smooth, then for any $x$ and $p$,
\[
f(x+p)\le f(x)+\nabla f(x)^{\top}p+\frac{L}{2}\|p\|^{2}.
\tag{6}
\]
\end{lemma}
\begin{proof}
Standard from Lipschitz continuity of the gradient. \qedhere
\end{proof}
\begin{lemma}[Cauchy Decrease]\label{lem:cauchy}
Let $g=\nabla f(x_k)$ and let $p_{\mathrm{c}}$ be the Cauchy point
\[
p_{\mathrm{c}}=-\frac{\|g\|}{\norm{B_k g}}\;g\;\text{ if }\;\norm{g}\neq0,
\]
clipped to the boundary $\|p_{\mathrm{c}}\|\le\Delta_k$. Then
\[
m_k(0)-m_k(p_{\mathrm{c}})
\ge \frac12\min\Bigl\{\frac{\|g\|^{2}}{L},\;\Delta_k\|g\|\Bigr\}.
\tag{7}
\]
\end{lemma}
\begin{proof}
Because $0\preceq B_k\preceq L I$, we have $\|B_k g\|\le L\|g\|$. Substituting the definition of $p_{\mathrm{c}}$ yields (7). \qedhere
\end{proof}

\section*{Main Proof (Step‑by‑Step Derivation)}
We prove Theorem \ref{thm:main}. Throughout the proof, let $g_k:=\nabla f(x_k)$.

\begin{enumerate}
\item[Step 1.] \textbf{Model overestimation.}  
From Lemma \ref{lem:descent} and $B_k\preceq L I$ we have for any feasible $p$,
\[
f(x_k+p)\le f(x_k)+g_k^{\top}p+\frac{L}{2}\|p\|^{2}
               \le f(x_k)+g_k^{\top}p+\frac12 p^{\top}B_k p
               = m_k(p).
\tag{8}
\]
Thus $f(x_k)-f(x_k+p_k)\ge \operatorname{pred}_k\cdot\rho_k$ by definition of $\rho_k$.

\item[Step 2.] \textbf{Lower bound on predicted reduction.}  
Using the Cauchy decrease guarantee \eqref{4},
\[
\operatorname{pred}_k
= m_k(0)-m_k(p_k)
\;\ge\; \kappa_{\mathrm{c}}\,
\min\!\Bigl\{\frac{\|g_k\|^{2}}{L},\;\Delta_k\|g_k\|\Bigr\}.
\tag{9}
\]

\item[Step 3.] \textbf{Acceptance condition.}  
If $\rho_k\ge \eta_1$, the trial step is accepted and the actual reduction satisfies
\[
\operatorname{ared}_k
= f(x_k)-f(x_{k+1})
\ge \eta_1\,\operatorname{pred}_k.
\tag{10}
\]

\item[Step 4.] \textbf{Decrease of the objective.}  
Combining (9) and (10) yields for every successful iteration,
\[
f(x_{k+1})\le f(x_k)-\eta_1\kappa_{\mathrm{c}}
\min\!\Bigl\{\frac{\|g_k\|^{2}}{L},\;\Delta_k\|g_k\|\Bigr\}.
\tag{11}
\]

\item[Step 5.] \textbf{Bounding the radius from below.}  
Whenever a step is rejected ($\rho_k<\eta_1$) the radius is multiplied by $\gamma_{\text{dec}}<1$. Since $\Delta_k$ can only be reduced a finite number of times before a successful step occurs (otherwise $f$ would decrease without bound, contradicting convexity), there exists a constant $\Delta_{\min}>0$ such that
\[
\Delta_k\ge \Delta_{\min}\quad\forall k.
\tag{12}
\]

\item[Step 6.] \textbf{Deriving a recursive inequality on the optimality gap.}  
Define $F_k:=f(x_k)-f^\star\ge0$. From convexity,
\[
F_k\le g_k^{\top}(x_k-x^\star).
\tag{13}
\]
Using Cauchy–Schwarz and (13),
\[
\|g_k\|\;\ge\;\frac{F_k}{\|x_k-x^\star\|}.
\tag{14}
\]
Since $f$ is $L$‑smooth, the iterates remain in the level set $\{x: f(x)\le f(x_0)\}$, which is bounded; denote $D:=\max_k\|x_k-x^\star\|<\infty$.

\item[Step 7.] \textbf{Two cases for the minimum in (11).}
\begin{description}
\item[Case A:] $\frac{\|g_k\|^{2}}{L}\le \Delta_k\|g_k\|$.  
Then (11) gives
\[
F_{k+1}\le F_k-\eta_1\kappa_{\mathrm{c}}\frac{\|g_k\|^{2}}{L}
        \le F_k-\frac{\eta_1\kappa_{\mathrm{c}}}{L}\frac{F_k^{2}}{D^{2}}
        \quad\text{(by (14))}.
\tag{15}
\]
\item[Case B:] $\Delta_k\|g_k\|\le \frac{\|g_k\|^{2}}{L}$.  
Using $\Delta_k\ge\Delta_{\min}$,
\[
F_{k+1}\le F_k-\eta_1\kappa_{\mathrm{c}}\Delta_{\min}\|g_k\|
        \le F_k-\eta_1\kappa_{\mathrm{c}}\Delta_{\min}\frac{F_k}{D}
        =F_k\Bigl(1-\frac{\eta_1\kappa_{\mathrm{c}}\Delta_{\min}}{D}\Bigr).
\tag{16}
\]
\end{description}
Both (15) and (16) imply a decrease of the form
\[
F_{k+1}\le F_k-\frac{\alpha}{k+1}F_k,
\qquad\alpha:=\frac{\eta_1\kappa_{\mathrm{c}}}{L}\frac{1}{D^{2}}.
\tag{17}
\]

\item[Step 8.] \textbf{Unrolling the recursion.}  
Inequality (17) yields
\[
\frac{1}{F_{k+1}}-\frac{1}{F_k}
   \ge \frac{\alpha}{k+1},
\]
and telescoping from $0$ to $k$ gives
\[
\frac{1}{F_{k+1}} \ge \frac{1}{F_0}+\alpha\sum_{i=0}^{k}\frac{1}{i+1}
               \ge \frac{1}{F_0}+\alpha\ln(k+2).
\]
Inverting,
\[
F_{k+1}\le \frac{1}{\frac{1}{F_0}+\alpha\ln(k+2)}
        \le \frac{C}{k+1},
\]
with $C:=\max\{F_0,\;2/\alpha\}$, establishing \eqref{5}.

\item[Step 9.] \textbf{Gradient‑norm complexity.}  
From (15) we have $F_{k+1}\le F_k-\frac{\eta_1\kappa_{\mathrm{c}}}{L D^{2}}F_k^{2}$.  
If $F_k>\varepsilon D$, then $F_{k+1}\le F_k-\frac{\eta_1\kappa_{\mathrm{c}}}{L D^{2}}(\varepsilon D)^{2}
 =F_k-\frac{\eta_1\kappa_{\mathrm{c}}\varepsilon^{2}}{L}$.  
Thus after at most
\[
N\le \frac{L D^{2}}{\eta_1\kappa_{\mathrm{c}}\varepsilon^{2}}
\]
successful steps the condition $F_k\le\varepsilon D$ holds, which together with (14) yields $\|g_k\|\le\varepsilon$. This proves Theorem \ref{thm:grad}.
\end{enumerate}

\section*{Rate Derivation (With Constants)}
Collecting the constants appearing in the proof:
\[
\alpha = \frac{\eta_1\kappa_{\mathrm{c}}}{L D^{2}},\qquad
C = \max\Bigl\{F_0,\;\frac{2}{\alpha}\Bigr\}
   =\max\Bigl\{f(x_0)-f^\star,\;
      \frac{2L D^{2}}{\eta_1\kappa_{\mathrm{c}}}\Bigr\}.
\]
Since $D\le \Delta_0+ \|x_0-x^\star\|$, a coarse bound $D\le \Delta_0+\|x_0-x^\star\|$ yields the constant stated in Theorem \ref{thm:main}.

\section*{Special Cases}
\begin{itemize}
\item \textbf{Exact Newton step.}  If $B_k=\nabla^{2}f(x_k)$ and the subproblem is solved exactly, then $p_k$ is the Newton direction truncated to the trust region. The same analysis applies, and if $\Delta_k$ is eventually larger than the Newton step length, quadratic convergence is recovered locally.
\item \textbf{Gradient descent as a limit.}  Setting $B_k=0$ reduces $m_k(p)=f(x_k)+g_k^{\top}p$, and the solution of \eqref{TR} is $p_k=-\min\{\Delta_k,\|g_k\|^{-1}\}g_k$. With a fixed radius $\Delta_k=\eta/L$, the method coincides with the classic $L$‑smooth gradient descent and the $O(1/k)$ rate follows directly.
\end{itemize}

\end{document}